% Part: sets-functions-relations
% Chapter: arithmetization
% Section: cauchy

\documentclass[../../../include/open-logic-section]{subfiles}

\begin{document}

\olfileid{sfr}{arith}{cauchy}
\olsection{Bijlage: de reële getallen als Cauchyrijen}
In \olref[cuts]{sec} hebben we de reële getallen geconstrueerd als
Dedekindsneden. In deze paragraaf leggen we een alternatieve
constructie uit. Die berust op Cauchy's definitie van wat we nu een
Cauchyrij noemen. Het gebruik van deze definitie om de reële getallen
te \emph{construeren} is echter afkomstig van andere negentiende-eeuwse
auteurs, met name Weierstrass, Heine, M\'{e}ray en Cantor. (Voor een
helder historisch overzicht, zie \citeauthor{OConnorRobertson:RN}
\citeyear{OConnorRobertson:RN}.)
Voordat we de negentiende eeuw behandelen, staan we stil bij Simon
Stevin (1548--1620). Stevin zag dat we elk reëel getal kunnen opvatten
via zijn decimaalontwikkeling. Zelfs een irrationaal getal als
$\sqrt{2}$ heeft dus een nette decimaalontwikkeling, die als volgt
begint:
\[
	1.41421356237\ldots
\] 
Decimaalontwikkelingen zijn eenvoudig in de verzamelingenleer te
modelleren: beschouw ze als functies $d \colon \Nat \to \Nat$, waarbij
$d(n)$ het $n$-de cijfer achter de komma is dat ons interesseert. Er is
dan een kleine aanpassing nodig voor het deel van het reële getal vóór
de komma (hier alleen $1$). Ook is een tweede aanpassing nodig, in de
vorm van een equivalentierelatie, om bijvoorbeeld te waarborgen dat
$0.999\ldots = 1$. Op deze manier kan binnen de verzamelingenleer zonder
veel moeite een volledig strenge constructie van de reële getallen in
de trant van Stevin worden gegeven.
Stevin is hier niet ons onderwerp. (Zie voor meer over Stevin
\citealt{KatzKatz2012}.) Wel gebruiken we een nauw verwant idee. In
plaats van de decimaalontwikkeling van $\sqrt{2}$ rechtstreeks te
beschouwen, kunnen we een \emph{rij} steeds nauwkeurigere rationale
benaderingen van $\sqrt{2}$ nemen, via de steeds preciezere
decimaalbenaderingen:
\[
1, 1.4, 1.414, 1.4142, 1.41421,\ldots
\]
Het idee dat reële getallen via ``steeds betere benaderingen'' kunnen
worden beschouwd, vormt de basis voor een nieuwe reeks inzichten. Ze
lijken op de inzichten waarmee we $\Rat$ uit $\Int$ en $\Int$ uit
$\Nat$ hebben geconstrueerd. Het gaat om het volgende:
\begin{enumerate}
	\item Elk reëel getal kan als een mogelijk oneindige
	decimaalontwikkeling worden geschreven.
	\item De informatie in zo'n mogelijk oneindige decimaalontwikkeling
	kan evengoed in een rij rationale getallen worden vastgelegd.
	\item Een rij rationale getallen kan worden opgevat als een functie
	van $\Nat$ naar $\Rat$: laat $f(n)$ eenvoudig het $n$-de rationale
	getal in de rij zijn.
\end{enumerate}
Niet \emph{elke} functie van $\Nat$ naar $\Rat$ levert echter een reëel
getal op. Beschouw bijvoorbeeld de functie:
\[
	f(n) =\begin{cases}
			1 & \text{als }n\text{ oneven is}\\
			0 &\text{als }n\text{ even is}
		\end{cases}
\]
Het probleem is dat de rij $0,1,0,1,0,1,0,\ldots$ geen enkel reëel
getal lijkt te benaderen. Om alleen rijen te beschouwen die wel een
reëel getal benaderen, beperken we ons daarom tot rijen met een
\emph{limiet}.
In \olref[his][set][limits]{sec} zijn we het begrip limiet al
tegengekomen. We kunnen hier echter niet \emph{precies} dezelfde
definitie gebruiken. De uitdrukking ``$(\forall \epsilon>0)$'' bevatte
daar stilzwijgend een kwantificatie over de reële getallen, en we
beschouwden limieten van functies op de reële getallen. Die definitie
gebruiken zou dus betekenen dat we de reële getallen al aannemen,
terwijl we juist die getallen willen \emph{construeren}. Gelukkig
kunnen we met een nauw verwant limietbegrip werken.
\begin{defn}\ollabel{def:CauchySequence}
  Een functie $f: \Nat \to \Rat$ is een \emph{Cauchyrij} dan en slechts
  dan als voor elke positieve $\epsilon \in \Rat$ geldt dat $(\exists
  \ell \in \Nat)(\forall m, n > \ell)|f(m) - f(n)| < \epsilon$.
\end{defn}
Het algemene limietidee is hetzelfde als voorheen. Voor elk gewenst
nauwkeurigheidsniveau, gemeten door~$\epsilon$, is er een ``gebied''
waarin we moeten kijken: alle invoeren groter dan~$\ell$. Onze rij
$1$, $1.4$, $1.414$, $1.4142$, $1.41421$\ldots heeft duidelijk een
limiet. Wie $\sqrt{2}$ wil benaderen met een fout kleiner dan
$\nicefrac{1}{10^{n}}$, kan elke term na de $n$-de nemen.
Een voor de hand liggende gedachte is dan dat elk reëel getal gewoon
een Cauchyrij \emph{is}. Net als bij de constructies van $\Int$ en
$\Rat$ zou dat te na\"ief zijn: hetzelfde reële getal wordt door
verschillende Cauchyrijen aangegeven. Dat is eenvoudig te zien. Neem
een Cauchyrij~$f$ en definieer $g$ als precies dezelfde functie als~$f$,
behalve dat $g(0)\neq f(0)$. Na de eerste term zijn de twee rijen
overal gelijk. Uiteindelijk willen we daarom zeggen dat ze dezelfde
limiet hebben, in de betekenis van \olref{def:CauchySequence}, en dus
hetzelfde reële getal ``definiëren''. We moeten deze Cauchyrijen dus als
hetzelfde reële getal beschouwen.
Daarom moeten we opnieuw een equivalentierelatie op de Cauchyrijen
definiëren en reële getallen met equivalentieklassen van Cauchyrijen identificeren.
Eerst hebben we het begrip nodig van een functie die in de limiet naar
$0$ gaat. Voor elke functie $h : \Nat \to \Rat$ zeggen we dat
\emph{$h$ naar $0$ gaat} dan en slechts dan als voor elke positieve
$\epsilon \in \Rat$ geldt dat $(\exists \ell \in \Nat)(\forall n >
\ell)|h(n)| < \epsilon$.\footnote{Vergelijk dit met de definitie van
$\lim_{x \mathord{\rightarrow}\infty}f(x) = 0$ in
\olref[his][set][limits]{sec}.} Zijn $f$ en $g$ functies $\Nat \to
\Rat$, stel dan $(f-g)(n) = f(n) - g(n)$. Definieer nu:
\[
	f \Realequiv g \text{ dan en slechts dan als $(f-g)$ naar $0$ gaat}.
\]
We moeten controleren dat $\Realequiv$ een equivalentierelatie is, en
dat is inderdaad zo. Vervolgens kunnen we de reële getallen desgewenst
definiëren als de equivalentieklassen onder $\Realequiv$ van alle
Cauchyrijen van $\Nat \to \Rat$.
\begin{prob}
Stel $f(n) = 0$ voor elke $n$. Stel $g(n) = \frac{1}{(n+1)^2}$. Toon
aan dat beide Cauchyrijen zijn en dat beide functies inderdaad limiet
$0$ hebben, zodat ook $f \sim_\Real g$.
\end{prob}
Zoals gebruikelijk schrijven we nu $\equivrep{f}{\Realequiv}$ voor de
equivalentieklasse met $f$ als !!a{element}. Voor de leesbaarheid laten
we de index hierna echter weg en schrijven we alleen
$\equivrep{f}{}$. Verder spreken we af dat voor elke $q \in \Rat$
geldt: $q_{\Real} = \equivrep{c_{q}}{}$, waarbij $c_{q}$ de constante
functie is met $c_q(n) = q$ voor alle $n \in \Nat$. Vervolgens
definiëren we de elementaire relaties en bewerkingen op de reële
getallen, bijvoorbeeld:
\begin{align*}
	\equivrep{f}{} + \equivrep{g}{} &= 	\equivrep{(f + g)}{} \\
	\equivrep{f}{} \times 	\equivrep{g}{} &= \equivrep{(f \times g)}{} 
\end{align*}
waarbij $(f + g)(n) = f(n) + g(n)$ en $(f \times g)(n) = f(n) \times
g(n)$. Ook moeten we controleren dat $(f + g)$, $(f-g)$ en $(f\times
g)$ Cauchyrijen zijn wanneer $f$ en $g$ dat zijn. Dat is inderdaad zo;
we laten het bewijs aan de lezer.
Ten slotte definiëren we een ordening. We noemen $\equivrep{f}{}$
\emph{positief} dan en slechts dan als zowel $\equivrep{f}{}\neq
0_\Real$ als $(\exists \ell \in \Nat)(\forall n > \ell)0 < f(n)$
geldt. Vervolgens definiëren we $\equivrep{f}{} < \equivrep{g}{}$ dan
en slechts dan als $\equivrep{(g - f)}{}$ positief is. We moeten
controleren dat dit welgedefinieerd is, dus niet afhangt van de keuze
van de ``representerende'' functie uit de equivalentieklasse.
%\begin{align*}
%	\equivrep{f}{} \leq \equivrep{g}{} \text{ iff }&\equivrep{f}{} = \equivrep{g}{} \text{ or }\\
%	&(\exists \ell \in \Nat)(\forall n \in \Nat)(\ell < n \lif f(n) \leq g(n))
%\end{align*}
Daarna is vrij eenvoudig aan te tonen dat deze definities de juiste
algebraïsche eigenschappen opleveren:
\begin{thm}\ollabel{thm:cauchyorderedfield}
De Cauchyrijen vormen een geordend lichaam.
\end{thm}
\begin{proof}
Oefening.
\end{proof}
\begin{prob}
Bewijs dat de Cauchyrijen een geordend lichaam vormen.
\end{prob}
Het is moeilijker te bewijzen dat de aldus geconstrueerde reële
getallen de volledigheidseigenschap hebben. Daarom geven we dat bewijs.
\begin{thm}
Elke niet-lege verzameling Cauchyrijen met een bovengrens heeft een
kleinste bovengrens.
\end{thm}
\begin{proof}[Bewijsschets] Zij $S$ een willekeurige niet-lege
verzameling Cauchyrijen met een bovengrens. Er is dus een $p \in \Rat$
waarvoor $p_{\Real}$ een bovengrens van $S$ is. Zij $r \in S$; dan is er
een $q \in \Rat$ waarvoor $q_{\Real} < r$. Als $S$ een kleinste
bovengrens heeft, ligt die dus tussen $q_\Real$ en $p_\Real$, met de
grenzen inbegrepen.
We sluiten de kleinste bovengrens steeds verder in door haar tegelijk
van onderen en van boven te benaderen. Daartoe definiëren we twee
functies $f, g \colon \Nat \to \Rat$, met het doel dat $f$ de kleinste
bovengrens\ van bovenaf en $g$ haar van onderaf benadert. We beginnen
met:
\begin{align*}
	f(0) &= p \\
	g(0) &= q
\end{align*}
Stel vervolgens $a_n = \frac{f(n) + g(n)}{2}$ en
definieer:\footnote{Dit is een recursieve definitie. We hebben echter
\emph{nog} geen reden gegeven om aan te nemen dat recursieve definities
toelaatbaar zijn.}
\begin{align*}
	f(n+1) &=
	\begin{cases}
		a_n &\text{als }(\forall h \in S)\equivrep{h}{} \leq (a_n)_\Real\\
		f(n)&\text{anders}
	\end{cases}\\
	g(n+1) &=
	\begin{cases}
		a_n &\text{als }(\exists h \in S)\equivrep{h}{} \geq (a_n)_\Real\\
	 	g(n) &\text{anders}
	\end{cases}
\end{align*}
Zowel $f$ als $g$ is een Cauchyrij. (Dit is vrij eenvoudig te
controleren, maar we laten het als oefening.) De functie $(f-g)$ gaat
naar $0$, want het verschil tussen $f$ en $g$ is na elke stap ten
hoogste half zo groot. Dus $\equivrep{f}{} = \equivrep{g}{}$.
Eerst tonen we aan dat $\equivrep{f}{}$ een bovengrens van $S$ is,
dus:\ dat $(\forall h \in S)\equivrep{h}{} \leq \equivrep{f}{}$. (Daarbij
gebruiken we \olref{thm:cauchyorderedfield}.) Zij $h \in S$ en neem ter
tegenspraak aan dat $\equivrep{f}{} < \equivrep{h}{}$, zodat $0_\Real <
\equivrep{(h-f)}{}$. Omdat $f$ een monotoon dalende Cauchyrij is, is er
een $n \in \Nat$ waarvoor $\equivrep{(c_{f(n)} - f)}{} <
\equivrep{(h-f)}{}$. Dus:
\[
	(f(n))_\Real = \equivrep{c_{f(n)}}{} < \equivrep{f}{} + \equivrep{(h-f)}{} = \equivrep{h}{},
\]
in tegenspraak met het feit dat volgens de constructie
$\equivrep{h}{} \leq (f(n))_\Real$.
Vervolgens tonen we aan dat $\equivrep{f}{} = \equivrep{g}{}$ de
\emph{kleinste} bovengrens van $S$ is. Zij $j$ een willekeurige
Cauchyrij en neem aan dat $\equivrep{j}{} < \equivrep{g}{}$. Net als
hierboven volgt, omdat $g$ \emph{stijgend} is, dat er een $n \in \Nat$
bestaat waarvoor $\equivrep{j}{} < (g(n))_\Real$. Volgens de constructie
is er echter een $h \in S$ waarvoor $(g(n))_\Real \leq
\equivrep{h}{}$. Dus $\equivrep{j}{} < \equivrep{h}{}$, zodat
$\equivrep{j}{}$ geen bovengrens van $S$ is.
\end{proof}
\end{document}

